\documentclass[a4paper]{article}

%------------------------------------------------
% الحزم
%------------------------------------------------

\usepackage{enumitem}
\usepackage[margin=2.5cm]{geometry}
\usepackage{amsmath,amsfonts,amssymb}
\usepackage[charter]{mathdesign}
\usepackage{graphicx}
\usepackage{array}

\usepackage{fancyhdr}
\usepackage{polyglossia}

%------------------------------------------------
% اللغات والخطوط
%------------------------------------------------

\setdefaultlanguage[
    calendar=gregorian,
    numerals=maghrib
]{arabic}

\setotherlanguage{english}

% إذا لم تكن الخطوط مثبتة في الجهاز، تُؤخذ من المجلد ../fonts/
\IfFontExistsTF{Amiri}{%
    \newfontfamily\arabicfont[Script=Arabic, Scale=1.2]{Amiri}
}{%
    \newfontfamily\arabicfont[
        Script=Arabic, Scale=1.2,
        Path=../fonts/, Extension=.ttf,
        UprightFont=*-Regular, BoldFont=*-Bold,
        ItalicFont=*-Italic, BoldItalicFont=*-BoldItalic
    ]{Amiri}
}

\IfFontExistsTF{AlBattar}{%
    \newfontfamily\arabicfontsf[Script=Arabic, Scale=1.3]{AlBattar}
    % خط العناوين (يعمل حتى بعد استعمال \textenglish)
    \newfontfamily\khatt[Script=Arabic, Scale=1.3]{AlBattar}
}{%
    \newfontfamily\arabicfontsf[
        Script=Arabic, Scale=1.3, Path=../fonts/,
        BoldFont=ae_AlBattar.ttf
    ]{ae_AlBattar.ttf}
    % خط العناوين (يعمل حتى بعد استعمال \textenglish)
    \newfontfamily\khatt[
        Script=Arabic, Scale=1.3, Path=../fonts/,
        BoldFont=ae_AlBattar.ttf
    ]{ae_AlBattar.ttf}
}

%------------------------------------------------
% إعداد الصفحة
%------------------------------------------------

\renewcommand{\footrulewidth}{1pt}
\renewcommand{\headrulewidth}{1pt}

\rhead{\khatt المدرسة العليا للأساتذة ببوزريعة}
\lhead{\khatt ملحقة بومرداس}

\cfoot{}
\lfoot{\khatt الأستاذ: د. بن عبد الرزاق يوسف}

\parindent 0pt
\setlength{\headsep}{15pt}
\setlength{\itemsep}{3cm}
\arraycolsep=1.4pt

\pagestyle{fancy}

\setlist[enumerate,1]{label=\arabic*)}


%================================================
% بداية الوثيقة
%================================================

\begin{document}

%------------------------------------------------
% العنوان
%------------------------------------------------

\begin{center}

{\Large\khatt\textbf{سلسلة تمارين}}

\vspace{0.2cm}

{\Large\khatt\textbf{في الذكاء الاصطناعي }}

\end{center}

\rule{\textwidth}{1pt}

\begin{center}

{\khatt\textbf{المحور الأول: اكتشاف الذكاء الاصطناعي التوليدي}}


\end{center}

\rule{\textwidth}{1pt}


%================================================
% النشاط الأول
%================================================

\section*{\khatt\underline{النشاط الأول: اكتشاف أدوات الذكاء الاصطناعي التوليدي}}

أصبح الذكاء الاصطناعي التوليدي من الأدوات التي يمكن للطالب
استعمالها في التعلم والبحث وإنجاز العديد من المهام.

من أشهر أدوات الذكاء الاصطناعي التوليدي:

\begin{itemize}

    \item \textenglish{ChatGPT}

    \item \textenglish{Gemini}

    \item \textenglish{Claude}

    \item \textenglish{DeepSeek}

    \item \textenglish{Qwen}

    \item \textenglish{Perplexity}

\end{itemize}

سنستعمل الأدوات الستة لتقديم نفس الطلب، ثم نقارن بين
الإجابات التي تقدمها.


\begin{enumerate}

%------------------------------------------------

\item أدخل الطلب الآتي في كل واحدة من الأدوات الستة:

\begin{center}

\textbf{
اشرح لطالب جامعي مبتدئ ما هو الذكاء الاصطناعي التوليدي،
واذكر ثلاثة أمثلة لاستعماله في الحياة اليومية.
}

\end{center}

%------------------------------------------------

\item سجّل الإجابة التي حصلت عليها من كل واحدة
من الأدوات الستة.

%------------------------------------------------

\item قارن بين إجابات الأدوات الستة من خلال الجدول الآتي:

\begin{center}

\resizebox{\textwidth}{!}{%

\begin{tabular}{|
>{\centering\arraybackslash}m{3cm}|
>{\centering\arraybackslash}m{2.2cm}|
>{\centering\arraybackslash}m{2.2cm}|
>{\centering\arraybackslash}m{2.2cm}|
>{\centering\arraybackslash}m{2.2cm}|
>{\centering\arraybackslash}m{2.2cm}|
>{\centering\arraybackslash}m{2.2cm}|
}

\hline

\textbf{المعيار}
&
\textenglish{\textbf{ChatGPT}}
&
\textenglish{\textbf{Gemini}}
&
\textenglish{\textbf{Claude}}
&
\textenglish{\textbf{DeepSeek}}
&
\textenglish{\textbf{Qwen}}
&
\textenglish{\textbf{Perplexity}}
\\

\hline

وضوح الإجابة
&
&
&
&
&
&
\\
\hline

سهولة الفهم
&
&
&
&
&
&
\\
\hline

تنظيم المعلومات
&
&
&
&
&
&
\\
\hline

عدد الأمثلة
&
&
&
&
&
&
\\
\hline

طول الإجابة
&
&
&
&
&
&
\\
\hline

\end{tabular}

}

\end{center}

%------------------------------------------------


\item غيّر صياغة الطلب السابق، واستعمل الطلب الآتي:

\begin{center}

\textbf{
اشرح لي الذكاء الاصطناعي التوليدي كما لو كنت طالبا
جامعيا مبتدئا، واستعمل أمثلة بسيطة من الحياة اليومية.
}

\end{center}

%------------------------------------------------

\item قارن بين الإجابة الجديدة والإجابة السابقة.

ماذا تلاحظ؟


\item اذكر ثلاثة أنواع مختلفة من المحتوى يستطيع
الذكاء الاصطناعي التوليدي إنتاجها.

%------------------------------------------------

\item أعط مثالين لاستعمال الذكاء الاصطناعي التوليدي
يمكن أن يساعدا الطالب في دراسته الجامعية.



\end{enumerate}

\newpage
%================================================
% النشاط الثاني
%================================================

\section*{\khatt\underline{النشاط الثاني (واجب منزلي): التحقق من إجابات الذكاء الاصطناعي}}

\textbf{يُنجز هذا النشاط في المنزل ويُسلَّم في الحصة القادمة.}

لا يعني الحصول على إجابة من الذكاء الاصطناعي أنها صحيحة
بالضرورة، لذلك يجب على الطالب قراءة الإجابة وتحليلها
والتحقق من المعلومات المهمة.


\begin{enumerate}

%------------------------------------------------

\item اختر إحدى الأدوات الستة السابقة، ثم أدخل الطلب الآتي:

\begin{center}

\textbf{
اكتب فقرة قصيرة حول فوائد ومخاطر استعمال الذكاء الاصطناعي
التوليدي في التعليم الجامعي.
}

\end{center}

%------------------------------------------------

\item اقرأ الإجابة التي قدمتها الأداة بعناية.

%------------------------------------------------

\item استخرج من الإجابة:

\begin{itemize}

    \item فائدتين للذكاء الاصطناعي التوليدي.

    \item خطرين محتملين.

    \item معلومة واحدة ترى أنها تحتاج إلى التحقق.

\end{itemize}

%------------------------------------------------

\item اختر المعلومة التي تحتاج إلى التحقق منها،
ثم ابحث عنها باستعمال مصدر موثوق.

%------------------------------------------------

\item قارن بين المعلومات التي وجدتها في المصدر
وبين إجابة الذكاء الاصطناعي.

%------------------------------------------------

\item هل وجدت تطابقا بين إجابة الذكاء الاصطناعي
والمصدر الذي اعتمدت عليه؟

علّل إجابتك.

%------------------------------------------------

\item ماذا تستنتج من هذه المقارنة؟

%------------------------------------------------

\item أمام العبارات الآتية، اكتب
\textbf{صحيح} أو \textbf{خطأ}:

\begin{enumerate}

    \item كل المعلومات التي يقدمها الذكاء الاصطناعي
    صحيحة دائما.

    \item يمكن للذكاء الاصطناعي التوليدي إنتاج النصوص.

    \item يمكن لبعض أدوات الذكاء الاصطناعي إنتاج الصور.

    \item يجب التحقق من المعلومات المهمة التي يقدمها
    الذكاء الاصطناعي.

    \item يمكن أن تتغير الإجابة حسب طريقة صياغة الطلب.

    \item يمكن للطالب استعمال الذكاء الاصطناعي كأداة
    مساعدة في التعلم.

\end{enumerate}

\end{enumerate}


\end{document}